\documentclass[10pt,a4paper]{amsart}
\usepackage[margin=19mm]{geometry}
\usepackage{amsmath,amssymb,mathtools}
\usepackage[scr=rsfs,cal=euler]{mathalfa}
\usepackage{xfrac}
\usepackage[unicode,hidelinks,bookmarks=false]{hyperref}
\newcommand{\ie}{i.e.\ }
\newcommand{\cf}{cf.\ }
\newcommand{\resp}{respectively\ }
\newcommand{\Subsection}{\subsection}
\providecommand{\nobreakdash}{\nobreak\hbox{-}}
\setlength{\emergencystretch}{2em}
\multlinegap0pt
\usepackage{xcolor}
\usepackage{amssymb}
\usepackage{float}
\usepackage{tikz}
\usepackage{algorithm}\usepackage{algpseudocode}
\newtheorem{lemma}{Lemma}
\newcommand{\R}{\mathbb{R}}
\newcommand{\szprox}{\ensuremath{\widehat{\operatorname{zprox}}}}
\newcommand{\zprox}{\ensuremath{\operatorname{zprox}}}
\title{Convergence of zeroth-order proximal point algorithms in the high-temperature regime}
\author{Emanuele Naldi, Hippolyte Labarri\`ere, Cesare Molinari, Silvia Villa}
\date{Source: arXiv:2605.11929v1 (CC BY 4.0)}
\begin{document}
\maketitle

\noindent\emph{Source and licence.} Convergence of zeroth-order proximal point algorithms in the high-temperature regime. Version arXiv:2605.11929v1 (CC BY 4.0), distributed by arXiv under the Creative Commons Attribution 4.0 International licence, http://creativecommons.org/licenses/by/4.0/, as recorded in the arXiv licence metadata for this exact version and shown on the abstract page https://arxiv.org/abs/2605.11929v1. Full licence notice: \texttt{licenses/arxiv-2605.11929-CC-BY-4.0.txt}.\par\smallskip

\noindent\textit{Editorial scope.} Counted unit: the article's lemma lem:boundedness (source0.tex lines 1642--1669). The article's complete original proof of it is delivered with it (source0.tex lines 1671--1855, one \texttt{proof} environment of 185 lines). No mathematics is written, repaired or re-derived: the statement and the proof are the article's own lines, in the article's own notation, and no retained line is substituted or re-flowed.  No further source-local context is needed: every cross-reference of every retained line resolves inside the two delivered spans.  Nothing was omitted from any retained span, so the package carries no omission record.  No retained line contains a citation, so the wrapper carries no bibliography and no bibliography entry.  No retained line contains a figure environment, an image-inclusion command or drawing code, so no image is delivered and the package declares no asset dependency.  Editorial formatting only, in addition: an AMS \texttt{amsart} preamble that restates the article's own declarations and macros used by the retained lines, namely the environments \texttt{lemma} on separate counters, together with the standard packages those lines need.  The retained environments print in this document as Lemma 1 in order of appearance, whereas the article prints the same items with the numbering of its own full document.  The complete native source of the article as distributed in the arXiv e-print for this exact version is bundled unchanged as \texttt{sources/200391/source0.tex}.
\begin{lemma}\label{lem:boundedness}
    Let \(f:\mathbb{R}^d\to \R\) be continuous and lower bounded. Then, for any compact set \(K\subset\R^d\) and any \(\varepsilon>0\), there exists \(C_{K,\varepsilon}>0\) such that
    \[
    \forall x\in K,\quad
    \mathbb{E}\left[
    \left\|
    \szprox^{\delta,N}_{\lambda, f}(x)
    -
    \zprox^{\delta}_{\lambda, f}(x)
    \right\|^2
    \mathbf{1}_{\Omega_{x,\varepsilon}}
    \right]
    \le
    \frac{C_{K,\varepsilon}}{N},
    \]
    where
    \[
    \Omega_{x,\varepsilon}
    =
    \left\{
    \omega:
    \frac{1}{N}\sum_{i=1}^N\exp(-f(y_i)/\delta)\ge\varepsilon
    \right\},
    \]
    and \(\left(y_i\right)_{i=1}^N\) are the i.i.d. samples from
    \(\mathcal{N}(x,\lambda\delta I)\) generating
    \(\szprox^{\delta,N}_{\lambda, f}(x)\).
\end{lemma}

\begin{proof}
We start by introducing some notation. For i.i.d. samples
\(y_1,\ldots,y_N\) from \(\mathcal{N}(x,\lambda \delta I)\), we denote
\[
    \mathbf u_i
    =
    y_i\exp(-f(y_i)/\delta),
    \qquad
    v_i
    =
    \exp(-f(y_i)/\delta),
\]
\[
    \hat{\mathbf u}
    =
    \frac{1}{N}\sum_{i=1}^N\mathbf u_i,
    \qquad
    \hat v
    =
    \frac{1}{N}\sum_{i=1}^Nv_i,
\]
and
\[
\begin{gathered}
\mathbf u
=
(2\pi\lambda\delta)^{-\frac{d}{2}}
\int_{\mathbb{R}^d}
y\,
\exp\!\left(-\frac{f(y)}{\delta}\right)
\exp\!\left(-\frac{\|x-y\|^2}{2\lambda\delta}\right)\,dy,
\\
v
=
(2\pi\lambda\delta)^{-\frac{d}{2}}
\int_{\mathbb{R}^d}
\exp\!\left(-\frac{f(y)}{\delta}\right)
\exp\!\left(-\frac{\|x-y\|^2}{2\lambda\delta}\right)\,dy.
\end{gathered}
\]
Then
\[
    \mathbb{E}[\hat{\mathbf u}]=\mathbf u,
    \qquad
    \mathbb{E}[\hat v]=v,
\]
and, for a given set of samples,
\[
    \szprox^{\delta,N}_{\lambda, f}(x)
    =
    \frac{\hat{\mathbf u}}{\hat v},
    \qquad
    \zprox^{\delta}_{\lambda, f}(x)
    =
    \frac{\mathbf u}{v}.
\]

First note that \(\mathbf u\) and \(v\) are continuous functions of \(x\), and consequently bounded on any compact set. In addition, since \(f\) is finite-valued, \(v(x)>0\) for every \(x\), and hence \(v\) is bounded away from zero on \(K\). Furthermore, since \(f\) is bounded below, for every \(p\ge 1\),
\[
\sup_{x\in K}\mathbb{E}\|\mathbf u_1\|^p<+\infty,
\qquad
\sup_{x\in K}\mathbb{E}|v_1|^p<+\infty .
\]
Indeed, if \(f\ge \inf f\), then
\[
\begin{aligned}
&\int_{\mathbb{R}^d}(1+\|y\|^p)
\exp\left(-\frac{p f(y)}{\delta}\right)
\exp\!\left(-\frac{\|x-y\|^2}{2\lambda\delta}\right)dy \\
&\qquad\le
\exp\left(-\frac{p\inf f}{\delta}\right)
\int_{\mathbb{R}^d}(1+\|y\|^p)
\exp\!\left(-\frac{\|x-y\|^2}{2\lambda\delta}\right)dy,
\end{aligned}
\]
and the last quantity is uniformly bounded for \(x\in K\).

On the event \(\Omega_{x,\varepsilon}=\{\hat v\ge \varepsilon\}\), the algebraic expansion
\[
\frac{\hat{\mathbf u}}{\hat v}-\frac{\mathbf u}{v}
=
\frac{\hat{\mathbf u}-\mathbf u}{v}
-
\frac{\mathbf u(\hat v-v)}{v^2}
-
\frac{(\hat{\mathbf u}-\mathbf u)(\hat v-v)}{v\hat v}
+
\frac{\mathbf u(\hat v-v)^2}{v^2\hat v}
\]
gives the pointwise bound
\[
\left\|
\szprox_{\lambda,f}^{\delta,N}(x)
-
\zprox_{\lambda,f}^{\delta}(x)
\right\|
\le A_1+A_2+A_3+A_4,
\]
where
\[
A_1:=\frac{\|\hat{\mathbf u}-\mathbf u\|}{v},
\quad
A_2:=\frac{\|\mathbf u\|\,|\hat v-v|}{v^2}, \quad A_3:=
\frac{\|\hat{\mathbf u}-\mathbf u\|\,|\hat v-v|}{v\varepsilon},
\quad
A_4:=
\frac{\|\mathbf u\|\,|\hat v-v|^2}{v^2\varepsilon}.
\]
Therefore, using \((\sum_{j=1}^4 A_j)^2\le 4\sum_{j=1}^4 A_j^2\), we obtain
\[
\begin{aligned}
\mathbb{E}\left[
\left\|
\szprox_{\lambda,f}^{\delta,N}(x)
-
\zprox_{\lambda,f}^{\delta}(x)
\right\|^2
\mathbf 1_{\Omega_{x,\varepsilon}}
\right] &\le
\frac{4}{v^2}\mathbb{E}\|\hat{\mathbf u}-\mathbf u\|^2
+
\frac{4\|\mathbf u\|^2}{v^4}\mathbb{E}|\hat v-v|^2 \\
&\quad+
\frac{4}{v^2\varepsilon^2}
\mathbb{E}\left[
\|\hat{\mathbf u}-\mathbf u\|^2|\hat v-v|^2
\right]
+
\frac{4\|\mathbf u\|^2}{v^4\varepsilon^2}
\mathbb{E}|\hat v-v|^4 .
\end{aligned}
\]
Since \(\hat{\mathbf u}\) and \(\hat v\) are empirical averages of i.i.d. random variables,
\[
\mathbb{E}\|\hat{\mathbf u}-\mathbf u\|^2
=
\frac{1}{N}\mathbb{E}\|\mathbf u_1-\mathbf u\|^2,
\qquad
\mathbb{E}|\hat v-v|^2
=
\frac{1}{N}\mathbb{E}|v_1-v|^2.
\]
Moreover, the fourth moments of \(\mathbf u_1\) and \(v_1\) are finite and uniformly bounded for \(x\in K\). Since the variables \(\mathbf u_i-\mathbf u\) and \(v_i-v\) are centered and i.i.d., the elementary estimate
\[
\mathbb{E}\left\|\frac1N\sum_{i=1}^N \xi_i\right\|^4
\le
\frac{C}{N^2}\mathbb{E}\|\xi_1\|^4
\]
for centered i.i.d. variables gives
\[
\mathbb{E}\|\hat{\mathbf u}-\mathbf u\|^4
\le
\frac{C_K}{N^2},
\qquad
\mathbb{E}|\hat v-v|^4
\le
\frac{C_K}{N^2}.
\]
By Cauchy--Schwarz,
\[
\mathbb{E}\left[
\|\hat{\mathbf u}-\mathbf u\|^2|\hat v-v|^2
\right]
\le
\sqrt{
\mathbb{E}\|\hat{\mathbf u}-\mathbf u\|^4
\mathbb{E}|\hat v-v|^4
}
\le
\frac{C_K}{N^2}.
\]
Since \(v\) is bounded away from zero on \(K\), and \(\mathbf u\) is bounded on \(K\), all prefactors are uniformly bounded. Therefore, there exists \(C_{K,\varepsilon}>0\) such that, for every \(x\in K\),
\[
\mathbb{E}\left[
\left\|
\szprox_{\lambda,f}^{\delta,N}(x)
-
\zprox_{\lambda,f}^{\delta}(x)
\right\|^2
\mathbf 1_{\Omega_{x,\varepsilon}}
\right]
\le
\frac{C_{K,\varepsilon}}{N}.
\]
\end{proof}

\end{document}
